% GENERATED FILE. Edit canonical lessons, then run npm run generate:books.
% canonical-source-hash: fnv1a64:8f4a4cc4999e6d16
% canonical-lessons: RU-C101-people, RU-C101-jumper, RU-C101-tie, RU-C101-gloves, RU-C101-scarf

\chapter{People, Jumper, Tie, Gloves, Scarf}
\label{ch:ru-people-jumper-tie-gloves-scarf}
\hlchaptermodality{101}

\begin{chapteropening}
\textbf{By the end of this chapter:} \emph{I can say people, a jumper, a tie, gloves and a scarf.}

Complete the last lesson of chapter 101: I can say people, a jumper, a tie, gloves and a scarf.
\end{chapteropening}

\section[\texorpdfstring{\ru{люди} (lyúdi)}{люди (lyúdi)}]{\ru{люди} (lyúdi) — people}

\label{lesson:RU-C101-people}



\begin{quote}
Before the new one: say the Russian for a calendar, then the Russian for a rucksack.
\end{quote}

\subsection*{You'll want to know: \ru{люди}}
\textbf{\ru{люди}} (\emph{lyúdi}) — \textquotedblleft{}people\textquotedblright{}.

\begin{cousinweb}
\textbf{\ru{люди}} $\leftarrow$ Proto-Slavic \emph{*ljudĭje}, \textquotedblleft{}people\textquotedblright{}, the plural of old \emph{\ru{люд}}, \textquotedblleft{}folk\textquotedblright{} — inherited. Vasmer relates it to Lithuanian \emph{liáudis}, \textquotedblleft{}the people\textquotedblright{}, Old High German \emph{liut}, \textquotedblleft{}people\textquotedblright{}, and Greek \emph{eleútheros}, \textquotedblleft{}free\textquotedblright{}, from a root \textquotedblleft{}to grow\textquotedblright{}.
\end{cousinweb}

\begin{grammarlens}[title={What you've built}]
One more word for this chapter.
\end{grammarlens}

\subsection*{Your turn}
\begin{itemize}
\raggedright
  \item \emph{Say it:} \emph{lyúdi}
  \item \emph{Say it:} \emph{lyúdi}, once more
  \item \emph{Say it:} say \emph{lyúdi}
  \item \emph{Recall:} say the Russian for a calendar, then the Russian for a rucksack, then say \emph{lyúdi} again
\end{itemize}

\begin{tcolorbox}[breakable,colback=teal!4,colframe=teal!35!black,title={Before you move on}]
What does \ru{люди} mean? (\textquotedblleft{}People\textquotedblright{}.) Say it once more.
\end{tcolorbox}
\section[\texorpdfstring{\ru{свитер} (svíter)}{свитер (svíter)}]{\ru{свитер} (svíter) — a jumper}

\label{lesson:RU-C101-jumper}



\begin{quote}
Before the new one: say the Russian for a rucksack, then the Russian for people.
\end{quote}

\subsection*{You'll want to know: \ru{свитер}}
\textbf{\ru{свитер}} (\emph{svíter}) — \textquotedblleft{}a jumper\textquotedblright{}.

\begin{cousinweb}
\textbf{\ru{свитер}} comes from English \emph{sweater}, from \emph{sweat}: first a garment for working up a sweat.
\end{cousinweb}

\begin{grammarlens}[title={What you've built}]
One more word for this chapter.
\end{grammarlens}

\subsection*{Your turn}
\begin{itemize}
\raggedright
  \item \emph{Say it:} \emph{svíter}
  \item \emph{Say it:} \emph{svíter}, once more
  \item \emph{Say it:} say \emph{svíter}
  \item \emph{Recall:} say the Russian for a rucksack, then the Russian for people, then say \emph{svíter} again
\end{itemize}

\begin{tcolorbox}[breakable,colback=teal!4,colframe=teal!35!black,title={Before you move on}]
What does \ru{свитер} mean? (\textquotedblleft{}A jumper\textquotedblright{}.) Say it once more.
\end{tcolorbox}
\section[\texorpdfstring{\ru{галстук} (gálstuk)}{галстук (gálstuk)}]{\ru{галстук} (gálstuk) — a tie}

\label{lesson:RU-C101-tie}



\begin{quote}
Before the new one: say the Russian for people, then the Russian for a jumper.
\end{quote}

\subsection*{You'll want to know: \ru{галстук}}
\textbf{\ru{галстук}} (\emph{gálstuk}) — \textquotedblleft{}a tie\textquotedblright{}.

\begin{cousinweb}
\textbf{\ru{галстук}} — first \emph{\ru{галздуки}}, in 1705, from Dutch \emph{halsdoek}, \textquotedblleft{}neck cloth\textquotedblright{}, later reshaped after German \emph{Halstuch} (Vasmer).
\end{cousinweb}

\begin{grammarlens}[title={What you've built}]
One more word for this chapter.
\end{grammarlens}

\subsection*{Your turn}
\begin{itemize}
\raggedright
  \item \emph{Say it:} \emph{gálstuk}
  \item \emph{Say it:} \emph{gálstuk}, once more
  \item \emph{Say it:} say \emph{gálstuk}
  \item \emph{Recall:} say the Russian for people, then the Russian for a jumper, then say \emph{gálstuk} again
\end{itemize}

\begin{tcolorbox}[breakable,colback=teal!4,colframe=teal!35!black,title={Before you move on}]
What does \ru{галстук} mean? (\textquotedblleft{}A tie\textquotedblright{}.) Say it once more.
\end{tcolorbox}
\section[\texorpdfstring{\ru{перчатки} (perchátki)}{перчатки (perchátki)}]{\ru{перчатки} (perchátki) — gloves}

\label{lesson:RU-C101-gloves}



\begin{quote}
Before the new one: say the Russian for a jumper, then the Russian for a tie.
\end{quote}

\subsection*{You'll want to know: \ru{перчатки}}
\textbf{\ru{перчатки}} (\emph{perchátki}) — \textquotedblleft{}gloves\textquotedblright{}.

\begin{cousinweb}
\textbf{\ru{перчатки}} comes from \emph{\ru{перст}}, an old word for \textquotedblleft{}finger\textquotedblright{} (Vasmer): first \emph{*pĭrstŭčatŭ}, \textquotedblleft{}fingered\textquotedblright{}, as in Old Russian \emph{\ru{пьрстатица}}, \textquotedblleft{}a glove\textquotedblright{}. Vasmer relates \emph{\ru{перст}} to Lithuanian \emph{pirštas}, \textquotedblleft{}finger\textquotedblright{}.
\end{cousinweb}

\begin{grammarlens}[title={What you've built}]
One more word for this chapter.
\end{grammarlens}

\subsection*{Your turn}
\begin{itemize}
\raggedright
  \item \emph{Say it:} \emph{perchátki}
  \item \emph{Say it:} \emph{perchátki}, once more
  \item \emph{Say it:} say \emph{perchátki}
  \item \emph{Recall:} say the Russian for a jumper, then the Russian for a tie, then say \emph{perchátki} again
\end{itemize}

\begin{tcolorbox}[breakable,colback=teal!4,colframe=teal!35!black,title={Before you move on}]
What does \ru{перчатки} mean? (\textquotedblleft{}Gloves\textquotedblright{}.) Say it once more.
\end{tcolorbox}
\section[\texorpdfstring{\ru{шарф} (sharf)}{шарф (sharf)}]{\ru{шарф} (sharf) — a scarf}

\label{lesson:RU-C101-scarf}



\begin{quote}
Before the new one: say the Russian for a tie, then the Russian for gloves.
\end{quote}

\subsection*{You'll want to know: \ru{шарф}}
\textbf{\ru{шарф}} (\emph{sharf}) — \textquotedblleft{}a scarf\textquotedblright{}.

\begin{cousinweb}
\textbf{\ru{шарф}} probably came through Polish \emph{szarfa}, \textquotedblleft{}a military sash\textquotedblright{}, from German \emph{Schärpe}, \textquotedblleft{}a sash\textquotedblright{}, which is traced to French \emph{écharpe}, \textquotedblleft{}an arm sling\textquotedblright{} (Vasmer).
\end{cousinweb}

\begin{grammarlens}[title={What you've built}]
One more word for this chapter.
\end{grammarlens}

\subsection*{Your turn}
\begin{itemize}
\raggedright
  \item \emph{Say it:} \emph{sharf}
  \item \emph{Say it:} \emph{sharf}, once more
  \item \emph{Say it:} say \emph{sharf}
  \item \emph{Recall:} say the Russian for a tie, then the Russian for gloves, then say \emph{sharf} again
\end{itemize}

\begin{tcolorbox}[breakable,colback=teal!4,colframe=teal!35!black,title={Before you move on}]
What does \ru{шарф} mean? (\textquotedblleft{}A scarf\textquotedblright{}.) Say it once more.
\end{tcolorbox}
